\PlanPage{11}{Section 4, Experiments: 7 pages, pp.~11--17. Section 4.1: 2 pages, pp.~11--12.}
\section{Experiments}\label{sec:evaluation}
\subsection{Experimental Setup}
\subsubsection{Datasets, Labels, and Temporal Splits}
\Budget{0.75 page on p.~11}
\DraftText{Identify the selected datasets (xxx), captures, protocol coverage,
attack labels, and target-flow labeling procedure. Provide raw-packet access,
timestamps, and endpoint information needed by both semantic layers. Specify
train/validation/test boundaries, host/session separation, and context reset
rules. Prevent overlapping communication history from leaking across splits.
State the observation cutoff used during evaluation.}
\DraftTable{4.0cm}{Datasets and evaluation splits}{Dataset/version / capture dates / flows / attack prevalence / protocols / available fields / split rule.\\
All dataset choices and counts: xxx.}
\DraftText{Define the cross-flow attack subset using independent annotations
or trace evidence. State capture-level or host-time-block units used for
uncertainty estimates and repeated runs.}
\subsubsection{Baselines and Comparison Conditions}
\Budget{0.75 page: 0.25 on p.~11 and 0.50 on p.~12}
\DraftText{Specify the same-task comparison groups: fast detector, suitable
trained traffic models, applicable cross-flow methods, Jev-IDS-style single-flow
decisions, and generative structured-output detectors. Final methods: xxx.
Define compatible prediction units before comparing them.}

\PlanPage{12}{Section 4.1 continued: baseline conditions 0.50 page; metrics 0.50 page.}
\DraftText{For complete-method comparisons, describe each method's input,
training budget, and reproduced implementation. For controlled model
substitution, match input evidence, detection operating points, and concurrency.
Use validation data to choose operating points. Identify any flow/trace
adaptation and record its effect on the task.}
\DraftTable{3.1cm}{Baseline implementations and controlled comparisons}{Method / prediction unit / available evidence / training or examples / model version / calibration / timed stages.}
\subsubsection{Metrics and Measurement Protocol}
\Budget{0.50 page on p.~12}
\DraftText{Report precision, recall, F1, FPR, and PR-AUC for the complete test
stream, with per-attack results. Report actual semantic requests per flow,
input tokens, cost, throughput, and p50/p95 latency. Include parsing, association,
context construction, API/network time, integration, and retries. Distinguish
observation waiting time from processing latency and per-call latency. State
hardware, concurrency, repetitions, and the uncertainty calculation.}
\DraftBlock{3.8cm}{Metric definitions and measurement boundaries.\\
Use consistent denominators and include all input flows in system-level results.}

\PlanPage{13}{Section 4.2: 1 page, p.~13. Main whole-system evidence.}
\subsection{End-to-End Detection and Same-Task Comparison}\label{sec:e2e}
\Budget{1 page}
\DraftText{Lead with full-stream target-flow results for the complete
framework and same-task methods. Quantify precision and recall changes relative
to the fast detector, then analyze attacks and protocols where related-flow
evidence explains the change. All measured values remain xxx until verified.}
\DraftTable{5.0cm}{End-to-end detection on the full test stream}{Dataset / method / precision / recall / F1 / FPR / PR-AUC / interval.\\[6pt]
Fast-only, selected same-task baselines, single-flow Jev, and the complete framework.\\[6pt]
Results: xxx.}
\DraftBlock{4.2cm}{Per-attack results and evidence-based interpretation.\\
Report the full evaluation population and all selected baselines.}

\PlanPage{14}{Section 4.3: 1 page, p.~14. Evidence for the semantic representation.}
\subsection{Semantic Representation Ablations}\label{sec:semantic-ablation}
\Budget{1 page}
\DraftText{Compare the full representation with removal of message-level
operation alignment, cross-flow context, deterministic statistics, and template
folding. Add equal-length random neighbors and perturb endpoint associations
or event order. Keep the target-flow evidence, Jev model, operating point, and
context budget controlled, and describe how each replacement is constructed.}
\DraftTable{3.8cm}{Semantic-module ablations}{Variant / precision / recall / FPR / input tokens / construction time.\\[6pt]
Full; raw-field input; target-only; no statistics; no folding; perturbed relations.\\
Results: xxx.}
\DraftFigure{3.3cm}{Effect of related-flow evidence under a matched context budget}{Real related flows versus random neighbors, endpoint perturbation, and order perturbation.\\
Cross-flow attack subset and uncertainty intervals: xxx.}
\DraftText{Explain which information each component contributes, including
the quality--budget tradeoff where a component compresses evidence.}

\PlanPage{15}{Section 4.4: 1 page, p.~15. Evidence for the two decision paths.}
\subsection{Decision-Path Ablations}\label{sec:path-ablation}
\Budget{1 page}
\DraftText{Evaluate fast-only, fast plus verification, fast plus discovery,
and the complete two-path framework on the same full test stream. Discovery
remains independently eligible across all flows. Report candidate counts,
overlap, actual calls, and the final detection metrics for every variant.}
\DraftTable{4.8cm}{Path contributions and error transitions}{Variant / precision / recall / corrected FP / lost TP / newly detected TP / added FP / calls.\\[6pt]
Fast-only; +verification; +discovery; both.\\
Results: xxx.}
\DraftBlock{4.3cm}{Interpretation of precision and recall changes.\\
Separate the effects of candidate selection, semantic reassessment, and final integration.\\
Include errors introduced as well as errors corrected.}

\PlanPage{16}{Section 4.5: 1 page, p.~16. Model substitution and whole-system cost.}
\subsection{Jev Comparison and End-to-End Efficiency}\label{sec:efficiency}
\Budget{1 page}
\DraftText{Replace Jev with the selected generative structured-output
decision model using matched evidence, detection operating points, and
concurrency. Record the model/version, input budget, decoding/output constraints,
and measurement period. Compare decision quality as well as latency and cost.}
\DraftFigure{4.0cm}{Detection quality and processing cost}{Quality versus p95 latency or throughput; annotate model, concurrency, and cost.\\[6pt]
Paired decision-model comparison and complete-pipeline measurements: xxx.}
\DraftTable{3.5cm}{End-to-end processing breakdown}{Parsing / association / routing / context building / API and network / integration.\\[6pt]
Throughput; p50/p95; calls per million flows; tokens; retries; cost. Results: xxx.}
\DraftText{Explain the separate effects of bounded contexts, typed decisions,
and invocation control. Relate efficiency claims to the measured complete
pipeline and report the time/cost of the actual API service configuration.}

\PlanPage{17}{Section 4.6: 1 page, p.~17. Budget sensitivity 0.60; traceable case 0.40.}
\subsection{Budget Sensitivity and Evidence Case Study}
\subsubsection{Context and Routing Budgets}
\Budget{0.60 page}
\DraftText{Vary the supported context window, token budget, and routing
budget. Show quality and cost jointly, and explain the operating point selected
on validation data. Separate changes in available evidence from changes in
decision-model execution.}
\DraftFigure{3.6cm}{Sensitivity to context and routing budgets}{Budget sweep: recall/FPR versus calls, tokens, throughput, and p95 latency.\\
Validation-selected operating point and test results: xxx.}
\subsubsection{Traceable Cross-Flow Evidence}
\Budget{0.40 page}
\DraftText{Follow one labeled target flow from messages and related
communications through the semantic context and decision paths to the final
output. Show source packet/flow IDs, time order, response changes, and the
ground-truth label. Use the case to explain the mechanism behind the aggregate
results.}
\DraftBlock{3.0cm}{Trace timeline, operational facts, source IDs, and target-flow decision.\\
Case and measured outputs: xxx.}
